%% kapa: exam=1 round=32 year=2021 period=1 form=A subject=civil_law subject_ko=민법 seq=12 number=12
%% source: Q-Net 감정평가사 기출문제 (cst00309) artlSeq=5208397 file=extracted/1008/5208397/A형/2021년도 제32회 감정평가사 1차 1교시 A형.pdf
%% conversion: kordoc 4.24.51 -> tools/gen1.py (auto); review=proofread-sonnet; flags=none
\begin{question}{12}
착오에 의한 의사표시에 관한 설명으로 옳지 않은 것은? (다툼이 있으면 판례에 따름)
\begin{choices}
  \item 토지매매에 있어서 특별한 사정이 없는 한, 매수인이 측량을 통하여 매매목적물이 지적도상의 그것과 정확히 일치하는지 확인하지 않은 경우 중대한 과실이 인정된다.
  \item 상대방이 표의자의 진의에 동의한 경우 표의자는 착오를 이유로 의사표시를 취소할 수 없다.
  \item 상대방에 의해 유발된 동기의 착오는 동기가 표시되지 않았더라도 중요부분의 착오가 될 수 있다.
  \item 상대방이 표의자의 착오를 알면서 이용한 경우에는 착오가 표의자의 중대한 과실로 인한 것이더라도 표의자는 착오에 의한 의사표시를 취소할 수 있다.
  \item 제3자의 기망행위에 의해 표시상의 착오에 빠진 경우에 사기가 아닌 착오를 이유로 의사표시를 취소할 수 있다.
\end{choices}
\end{question}
\answer{1}
